\section{Deployment and Runtime Configuration}\label{sec:deployment}

The local production-like stack is defined by Docker Compose.
The server image builds the C++ compiler and installs \texttt{motivoc} at \texttt{/usr/local/bin/motivoc}.
The client image builds the Next.js application and proxies API requests to the server inside the Compose network.
The nginx gateway terminates HTTPS with a local self-signed certificate and routes browser traffic to the client and
API routes to the server.
Docker is used here as a reproducibility tool: it keeps the compiler, server, and client runtime assumptions explicit
instead of relying on an ad hoc local machine setup~\cite{Docker_2026}.

\begin{figure}[htbp]
  \centering
  \begin{tikzpicture}[
      svc/.style={rounded corners=5pt, draw=motivoInk!65, very thick, minimum width=3.8cm, minimum height=1.1cm,
      align=center, font=\small\sffamily},
      arrow/.style={-Stealth, very thick, draw=motivoInk!75}
    ]
    \node[svc, fill=motivoSoft] (browser) {Browser\\\texttt{motivo-studio.local}};
    \node[svc, fill=motivoAmber!18, below=0.8cm of browser] (nginx) {nginx gateway\\ports 80 and 443};
    \node[svc, fill=motivoBlue!13, below left=0.9cm and 0.7cm of nginx] (client) {Next.js client\\port 3000};
    \node[svc, fill=motivoCyan!16, below right=0.9cm and 0.7cm of nginx] (server) {Express server\\port 3001};
    \node[svc, fill=motivoGreen!16, below=0.8cm of server] (compiler) {\texttt{motivoc}\\inside server image};

    \draw[arrow] (browser) -- (nginx);
    \draw[arrow] (nginx) -- node[left, font=\scriptsize\sffamily] {page assets} (client);
    \draw[arrow] (nginx) -- node[right, font=\scriptsize\sffamily] {\texttt{/api/*}} (server);
    \draw[arrow] (server) -- (compiler);
  \end{tikzpicture}
  \caption{Docker Compose deployment topology for Motivo Studio.}\label{fig:deployment}
\end{figure}

Configuration is intentionally small.
The server reads \texttt{PORT}, \texttt{HOSTNAME}, and \texttt{COMPILER\_BIN}.
The client uses \texttt{API\_PROXY\_URL} for API rewrites.
This is sufficient for the current deployment model, where the compiler is a local executable reachable by the API
container.
Future remote compilation would require a different trust and resource model, but that is outside the present thesis.
